\documentclass[11pt]{report}

\usepackage[a4paper,margin=25mm]{geometry}
\usepackage{amsmath}
\usepackage{amssymb}
\usepackage{bm}
\usepackage{booktabs}
\usepackage{longtable}
\usepackage{mathtools}
\usepackage{microtype}
\usepackage[hidelinks]{hyperref}

\hypersetup{
  pdftitle={Atomistic DWBA for Unit Cells in a Graded Interface},
  pdfauthor={orGUI design note},
  pdfsubject={Four-channel unit-cell DWBA and specular optical-reference subtraction}
}

\newcommand{\dd}{\mathop{}\!\mathrm{d}}
\newcommand{\ii}{\mathrm{i}}
\newcommand{\ee}{\mathrm{e}}
\newcommand{\kzero}{k_0}
\newcommand{\Qpar}{\bm{Q}_{\parallel}}
\newcommand{\rpar}{\bm{r}_{\parallel}}
\newcommand{\Cchan}{\mathcal C}
\newcommand{\Phiint}{\Phi}

\title{Atomistic DWBA for Unit Cells in a Graded Interface}
\author{orGUI design note}
\date{\today}

\begin{document}

\hypersetup{pageanchor=false}
\maketitle

\begin{abstract}
This note specifies the first-order distorted-wave Born approximation needed
to add finite atomic models to orGUI's graded optical reference.  It derives,
rather than quotes, both the matrix element and the factor that converts it
into a measurable amplitude: Born expansion, Green function of the stratified
reference built explicitly from two homogeneous solutions and their Wronskian,
lattice expansion of the index contrast, and lateral collapse.  The Wronskian
is shown to equal $2\ii k_{z,0}$ independently of the prepared stack, which is
the origin of every factor $1/k_z$ in this subject; the reciprocal exit field
then follows from a boundary-value problem, with no far-field approximation
over the illuminated footprint --- an approximation that is in fact invalid on
the specular rod.

The result is a single matrix element $F_{\bm h}$ per crystal truncation rod,
normalised to one reference lateral cell.  A generated \texttt{UnitCell},
whether it belongs to a \texttt{Film}, an \texttt{EpitaxyInterface}, or a
surface model, is evaluated in every optical medium that contains its atoms,
and contributes to $F_{\bm h}$ as the coherent sum of the four Renaud
incident/exit channels --- which are simply the two branches of each reference
field.  Because the prepared optical reference is laterally uniform, its
constant $\delta_m+\ii\beta_m$ density enters the $\bm h=\bm 0$ Fourier
coefficient and no other: it is part of $F_{\bm 0}$, not something subtracted
afterwards from an atomic sum, and it is absent from every genuine rod.  That
contribution is partitioned among unit cells without changing the
profile-level result, so signed occupancies, strain, lateral-area conversion,
and optical-profile simplification remain consistent.
\end{abstract}

\tableofcontents
\clearpage
\hypersetup{pageanchor=true}
\pagenumbering{arabic}

\chapter{Reference problem, conventions, and the master formula}
\label{chap:reference}

\section{Scope}

The unperturbed electric fields and the unperturbed specular reflection
amplitude are assumed to be available from the layered-field solver.  The
problem addressed here is the remaining atomistic matrix element for finite
unit cells in the graded interface region.  The terminal semi-infinite bulk is
included only to show that the new formula reduces to the already implemented
bulk result.

The derivation is electromagnetic throughout.  Following Renaud, Lazzari, and
Leroy, the incident field and the reciprocal field launched from the detector
are denoted by electric fields $\bm E_0$; no quantum-mechanical wavefunction
notation is used.  The medium is nonmagnetic, isotropic, linear, and reciprocal.
Magnetic and tensor-resonant scattering are outside the present scope.

This chapter derives, rather than quotes, both the matrix element and the
factor that converts it into a measurable rod amplitude.  The chain is: Born
expansion, Green function of the stratified reference built from two
homogeneous solutions and their Wronskian, lattice expansion of the index
contrast, and lateral collapse.  No far-field (Fraunhofer) approximation over
the illuminated footprint is used at any point; Sec.~\ref{sec:why-not-farfield}
explains why one must not be.

\section{Fourier, surface, and refractive-index conventions}

We use the crystallographic space-time convention
\begin{equation}
  \bm E(\bm r,t)\propto
  \exp\!\left[\ii\left(\omega t-\bm k\mathbin{\cdot}\bm r\right)\right],
  \qquad
  \kzero=\frac{2\pi}{\lambda},
  \qquad
  \bm Q=\bm k_f-\bm k_i,
  \label{eq:convention}
\end{equation}
and the Fourier transform
\begin{equation}
  \widetilde g(\bm Q)
  =\int g(\bm r)\exp(\ii\bm Q\mathbin{\cdot}\bm r)\,\dd^3r.
  \label{eq:fourier}
\end{equation}
The $z$ axis is the outward surface normal: vacuum is at large positive $z$
and the bulk extends toward negative $z$.  Momentum transfer is in
Angstrom$^{-1}$ and all real-space coordinates below are in Angstrom.

The prepared optical reference consists of planar media $m=0,\ldots,M$.  In
medium $m$,
\begin{equation}
  n_m=1-\delta_m-\ii\beta_m,
  \qquad \delta_m,\beta_m\in\mathbb R,
  \label{eq:index}
\end{equation}
and the medium occupies the interval
\begin{equation}
  I_m=(z_m^-,z_m^+).
\end{equation}
One or both endpoints may be infinite for a terminal medium.  A fixed
half-open convention must be used in code when an atom lies exactly on a
boundary.

orGUI uses the negative-root Renaud normal wavevector
\begin{equation}
  k_{z,m}=-\kzero
  \sqrt{n_m^2-n_0^2\cos^2\alpha},
  \qquad
  \operatorname{Re}k_{z,m}\leq0,
  \quad
  \operatorname{Im}k_{z,m}\geq0.
  \label{eq:renaud-kz}
\end{equation}
The incident and reciprocal exit solutions are calculated in the same frozen
array of complex refractive indices and at the same physical boundaries.

Equation~\eqref{eq:renaud-kz} fixes what ``upward'' means, and every sign in
this document follows from it.  Because $\operatorname{Re}k_{z,m}\leq0$, a
wave travelling towards $+z$ --- upward, out of the sample and towards the
detector --- carries $\exp(+\ii k_{z,m}z)$, while a wave travelling towards
$-z$, downward onto the sample, carries $\exp(-\ii k_{z,m}z)$.  This is the
opposite pairing from the one a positive-root convention would give, and it is
the single most common source of sign errors when comparing with the
literature.  The whole of Secs.~\ref{sec:born}--\ref{sec:collapse} is carried
out in Eq.~\eqref{eq:renaud-kz}'s convention, so no conversion is ever needed.

\section{Scattering-factor density and the index contrast}

For atom $a$, orGUI uses
\begin{equation}
  f_a(Q^2,E)=f^0_{a,\mathrm{WK}}(Q^2)+f'_a(E)+\ii f''_a(E),
  \label{eq:atomic-factor}
\end{equation}
where the Waasmaier--Kirfel representation is evaluated analytically at the
generally complex channel value $Q^2$.  Its zero-transfer value, including the
same anomalous terms, is also used to generate the optical profile.  Replacing
that value by an independent atomic-number approximation would break the
reference-density cancellation.  Equation~\eqref{eq:atomic-factor} is the only
quantity in this document called $f$.

Let $\rho_f(\bm r)$ be the complex scattering-factor density in electrons per
Angstrom$^3$.  To first order in the X-ray susceptibility,
\begin{equation}
  n^2(\bm r)\simeq
  1-\frac{4\pi r_e}{\kzero^2}\rho_f(\bm r),
  \label{eq:index-density}
\end{equation}
with $r_e$ the classical electron radius.  The piecewise-constant reference
density represented by Eq.~\eqref{eq:index} is therefore
\begin{equation}
  \boxed{
  \overline\rho_{f,m}
  =\frac{\kzero^2}{2\pi r_e}
   \left(\delta_m+\ii\beta_m\right).
  }
  \label{eq:reference-density}
\end{equation}
Write $n_0^2(z)=\sum_m n_m^2\,\mathbf 1_{I_m}(z)$ for the prepared, laterally
uniform reference and
\begin{equation}
  \delta n^2(\bm r)=n^2(\bm r)-n_0^2(z)
  =-\frac{4\pi r_e}{\kzero^2}\,\Delta\rho_f(\bm r),
  \qquad
  \Delta\rho_f(\bm r)
  =\rho_f(\bm r)-\sum_m\overline\rho_{f,m}\,\mathbf 1_{I_m}(z),
  \label{eq:density-contrast}
\end{equation}
for the index contrast that drives all scattering.  Equations~\eqref{eq:index}
and \eqref{eq:reference-density} enter the problem only through $n_0^2(z)$, and
Sec.~\ref{sec:collapse} shows exactly where they survive.

\section{The Born expansion}
\label{sec:born}

At the small scattering angles of interest the vectorial propagation equation
reduces to a scalar Helmholtz equation for the component of the field along
the analysed polarization (Renaud \S5.3.6).  Writing $E$ for that component,
\begin{equation}
  \left[\nabla^2+\kzero^2n^2(\bm r)\right]E(\bm r)=0,
  \qquad
  n^2(\bm r)=n_0^2(z)+\delta n^2(\bm r),
  \label{eq:helmholtz}
\end{equation}
so that
\begin{equation}
  \left[\nabla^2+\kzero^2n_0^2(z)\right]E(\bm r)
  =-\kzero^2\,\delta n^2(\bm r)\,E(\bm r).
  \label{eq:driven}
\end{equation}
Defining the Green function of the \emph{reference} problem by
\begin{equation}
  \left[\nabla^2+\kzero^2n_0^2(z)\right]G(\bm r,\bm r')=\delta(\bm r-\bm r'),
  \label{eq:green-def}
\end{equation}
Eq.~\eqref{eq:driven} is equivalent to the integral equation
\begin{equation}
  E(\bm r)=\bm E_0(\bm r)
  -\kzero^2\!\int\!\dd^3r'\,G(\bm r,\bm r')\,\delta n^2(\bm r')\,E(\bm r'),
  \label{eq:lippmann}
\end{equation}
whose iteration is the Born series.  Retaining the first order, with $\bm E_0$
and $G$ both those of the \emph{stratified} reference, gives the distorted-wave
Born approximation:
\begin{equation}
  \boxed{
  \delta E(\bm r)
  =-\kzero^2\!\int\!\dd^3r'\,G(\bm r,\bm r')\,\delta n^2(\bm r')\,
   \bm E_0(\bm r',\bm k_i).
  }
  \label{eq:dwba-master}
\end{equation}
Reflection, transmission, refraction, and repeated reflection at the planar
interfaces are already contained in $\bm E_0$ and in $G$; the perturbation is
scattered exactly once.

\paragraph{Sign remark.}  Renaud's Eqs.~(71) and (73) as printed carry
$+\kzero^2\delta n^2$ on the right-hand side.  With his Eq.~(70)'s definition
$\delta n^2=n^2-n_0^2$, adopted here in Eq.~\eqref{eq:density-contrast}, the
sign is as in Eqs.~\eqref{eq:driven} and \eqref{eq:lippmann}.  The overall sign
of the final result is verified independently against the Born--Fresnel limit
in Sec.~\ref{sec:born-check}.

\section{The Green function of the stratified reference}
\label{sec:green}

This section does the work that is usually skipped.  Nothing in it is
approximate.

\subsection{Separation into a one-dimensional problem}

Because $n_0^2$ depends only on $z$, Eq.~\eqref{eq:green-def} separates under a
lateral Fourier transform:
\begin{equation}
  G(\bm r,\bm r')
  =\frac{1}{(2\pi)^2}\int\dd^2q\;g_q(z,z')\,
   \ee^{-\ii\bm q\cdot(\rpar-\rpar')},
  \qquad
  \left[\frac{\dd^2}{\dd z^2}+\kzero^2n_0^2(z)-q^2\right]g_q(z,z')
  =\delta(z-z').
  \label{eq:separation}
\end{equation}
The lateral wavevector $\bm q$ enters the one-dimensional operator only through
the constant $-q^2$; the normal wavevector belonging to $\bm q$ in medium $m$
is Eq.~\eqref{eq:renaud-kz}'s $k_{z,m}$.

\subsection{Wronskian construction}

Equation~\eqref{eq:separation} is a one-dimensional Green-function problem for
an operator with no first-derivative term.  Let $u_-$ and $u_+$ solve the
homogeneous equation subject to one physical boundary condition each,
\begin{align}
  u_-(z):&\quad\text{no wave incoming from below, as }z\to-\infty,
  \label{eq:bc-minus}\\
  u_+(z):&\quad\text{no wave incoming from above, as }z\to+\infty,
  \label{eq:bc-plus}
\end{align}
each fixed up to one multiplicative constant.  Then
\begin{equation}
  g_q(z,z')=\frac{u_-(z_<)\,u_+(z_>)}{W},
  \qquad
  W=u_-u_+'-u_-'u_+,
  \label{eq:wronskian-construction}
\end{equation}
with $z_<=\min(z,z')$ and $z_>=\max(z,z')$.  The Wronskian is independent of
$z$,
\begin{equation}
  \frac{\dd W}{\dd z}
  =u_-u_+''-u_-''u_+
  =u_-\!\left(q^2-\kzero^2n_0^2\right)u_+
   -\left(q^2-\kzero^2n_0^2\right)u_-u_+=0,
  \label{eq:wronskian-constant}
\end{equation}
precisely because the operator has no $\dd/\dd z$ term, so it may be evaluated
wherever it is easiest.  The normalisation of $u_\pm$ cancels in
Eq.~\eqref{eq:wronskian-construction}.

\subsection{What the two solutions are physically}

Using the convention established after Eq.~\eqref{eq:renaud-kz}, in vacuum
($m=0$):
\begin{itemize}
  \item $u_+$ is fixed by Eq.~\eqref{eq:bc-plus}: above the whole structure it
        is purely upward-going,
        \begin{equation}
          u_+(z)=\ee^{+\ii k_{z,0}z}.
          \label{eq:u-plus}
        \end{equation}
        How it continues downward through the stack is never needed, because
        the observation point is always in vacuum.
  \item $u_-$ is fixed by Eq.~\eqref{eq:bc-minus}: deep in the terminal medium
        it carries only the transmitted, downward or evanescently decaying
        wave, which is the physical condition $A_M^-=0$.  That single condition,
        continued \emph{upward} through the stack, produces in vacuum a
        superposition of a downward and an upward component; normalising the
        downward one to unity,
        \begin{equation}
          u_-(z)=\ee^{-\ii k_{z,0}z}+r_0\,\ee^{+\ii k_{z,0}z}
          \qquad\text{(in vacuum)}.
          \label{eq:u-minus}
        \end{equation}
        This is the ordinary Fresnel/Parratt reference field for a plane wave
        incident from above --- incident plus reflected, with only a
        transmitted wave below --- and $r_0$ is the unperturbed specular
        reflection amplitude of Eq.~\eqref{eq:specular-reflectivity}.  The
        condition that \emph{defines} $u_-$ is imposed at $z\to-\infty$, but
        the solution it selects is the familiar one driven from above.
\end{itemize}
Since the exit wavevector $\bm k_f$ points away from the sample, $-\bm k_f$
points onto it, and $u_-$ evaluated for that direction is exactly the
reciprocal exit field,
\begin{equation}
  u_-(z')=\bm E_0(z',-\bm k_f).
  \label{eq:u-minus-identification}
\end{equation}
This is the content of the reciprocity theorem, obtained here by solving a
boundary-value problem rather than by propagating a field to a distant
detector.  It is why the reciprocal exit field is never complex conjugated:
it is a solution of the same homogeneous equation as the incident field, not
its adjoint.

\subsection{The Wronskian, and the Green function}

Evaluate $W$ in vacuum from Eqs.~\eqref{eq:u-plus}--\eqref{eq:u-minus}:
\begin{align}
  W&=\left[\ee^{-\ii k_{z,0}z}+r_0\ee^{\ii k_{z,0}z}\right]
     \left(\ii k_{z,0}\ee^{\ii k_{z,0}z}\right)
    -\left[-\ii k_{z,0}\ee^{-\ii k_{z,0}z}
       +\ii k_{z,0}r_0\ee^{\ii k_{z,0}z}\right]\ee^{\ii k_{z,0}z}\nonumber\\
   &=\ii k_{z,0}+\ii k_{z,0}r_0\ee^{2\ii k_{z,0}z}
     +\ii k_{z,0}-\ii k_{z,0}r_0\ee^{2\ii k_{z,0}z}
   \;=\;2\ii k_{z,0}.
  \label{eq:wronskian-value}
\end{align}
The $r_0$ terms cancel identically: \emph{the Wronskian does not depend on the
stack}, only on the vacuum $k_{z,0}$.  This one fact is the origin of every
factor $1/k_z$ in this subject, including the one that appears without
explanation in the reflectivity literature.

With the observation point above all sources, $z>z'$, so $z_<=z'$, $z_>=z$, and
Eqs.~\eqref{eq:wronskian-construction}, \eqref{eq:u-plus},
\eqref{eq:u-minus-identification} and \eqref{eq:wronskian-value} give
\begin{equation}
  \boxed{
  g_q(z,z')
  =\frac{\bm E_0(z',-\bm k_f)\,\ee^{\ii k_{z,0}z}}{2\ii k_{z,0}}
  =-\frac{\ii}{2k_{z,0}}\,\ee^{\ii k_{z,0}z}\,\bm E_0(z',-\bm k_f).
  }
  \label{eq:green-1d}
\end{equation}
For a homogeneous reference, $r_0=0$ and $u_-(z')=\ee^{-\ii k_{z,0}z'}$, so
Eq.~\eqref{eq:green-1d} collapses to
$g_q=-(\ii/2k_{z,0})\ee^{\ii k_{z,0}|z-z'|}$, whose derivative jump across
$z=z'$ is unity; substituting that into Eq.~\eqref{eq:separation} reproduces
the Weyl angular-spectrum representation of an outgoing spherical wave.  The
Weyl expansion is therefore not an extra ingredient but the vacuum special case
of Eq.~\eqref{eq:green-1d}.

\subsection{Why the far-field route is not available here}
\label{sec:why-not-farfield}

The textbook shortcut replaces $G$ by its far-field limit, expanding
$|\bm r-\bm r'|\simeq r-\bm k_f\cdot\bm r'/\kzero$ and pulling a single
spherical wave $\ee^{-\ii\kzero r}/r$ out of the integral.  That expansion
requires the detector distance $R$ to exceed the Fraunhofer distance of the
\emph{whole illuminated footprint}.  For a grazing-incidence footprint of
lateral size $L\sim1$~mm to 1~cm at $\lambda\sim1$~\AA, the neglected
quadratic phase is
\begin{equation}
  \delta(\text{phase})\sim\kzero\frac{L^2}{R}\sim10^{4}\text{--}10^{7}\gg1,
  \label{eq:fraunhofer-failure}
\end{equation}
so the far-field form of $G$ is never valid on the specular rod.
Equation~\eqref{eq:green-1d} avoids the issue entirely: it is exact at any $R$.

The polarization factor is a separate matter.  The same far-field substitution
is conventionally used to fix the exit polarization vector, but there it enters
undifferentiated and is \emph{not} multiplied by $\kzero$, so its error is only
$\sim L/R\sim10^{-2}$ rather than Eq.~\eqref{eq:fraunhofer-failure}.  Fixing the
polarization basis at the nominal $\bm k_f$, as
Eq.~\eqref{eq:polarization-contractions} does, therefore remains accurate even
where the phase treatment above is mandatory.

\section{Lattice expansion and the lateral collapse}
\label{sec:collapse}

\subsection{The index contrast along the surface lattice}

Parallel to the surface, $n^2(\bm r)$ is periodic on the two-dimensional
lattice of the reference lateral cell, of area $A_{\rm ref}$, while $n_0^2(z)$
is laterally uniform.  Their difference $\delta n^2$ is therefore
two-dimensional-periodic and has the Fourier series
\begin{equation}
  \delta n^2(\rpar,z)=\sum_{\bm h}\delta n^2_{\bm h}(z)\,
  \ee^{-\ii\bm h\cdot\rpar},
  \qquad
  \delta n^2_{\bm h}(z)
  =\frac{1}{A_{\rm ref}}\int_{\rm cell}\dd^2r_\parallel\;
   \delta n^2(\rpar,z)\,\ee^{\ii\bm h\cdot\rpar},
  \label{eq:lattice-expansion}
\end{equation}
where $\bm h$ runs over the reciprocal lattice of that cell, i.e.\ over the
crystal truncation rods $(h,k)$.  Because $n_0^2(z)$ is laterally uniform it
contributes to Eq.~\eqref{eq:lattice-expansion} at $\bm h=\bm 0$ and nowhere
else:
\begin{align}
  \delta n^2_{\bm h}(z)
  &=\frac{1}{A_{\rm ref}}\int_{\rm cell}\dd^2r_\parallel\;n^2(\rpar,z)\,
    \ee^{\ii\bm h\cdot\rpar},
  &&\bm h\neq\bm0,
  \label{eq:coeff-nonspec}\\
  \delta n^2_{\bm 0}(z)
  &=\frac{1}{A_{\rm ref}}\int_{\rm cell}\dd^2r_\parallel\;n^2(\rpar,z)
    \;-\;n_0^2(z),
  &&\bm h=\bm 0.
  \label{eq:coeff-spec}
\end{align}
Equations~\eqref{eq:coeff-nonspec}--\eqref{eq:coeff-spec} are the entire
content of the specular/non-specular distinction, and they are the reason this
document needs only one matrix element.  The $\delta_m,\beta_m$ of the prepared
profile are not something subtracted afterwards from an atomic sum: they are
part of the $\bm h=\bm 0$ Fourier coefficient, and part of nothing else.

\subsection{Collapse}

Insert Eqs.~\eqref{eq:separation}, \eqref{eq:green-1d} and
\eqref{eq:lattice-expansion} into Eq.~\eqref{eq:dwba-master}.  Stratification
lets the lateral phase of the incident field be factored out,
$\bm E_0(\bm r',\bm k_i)
=\ee^{-\ii\bm k_{i\parallel}\cdot\rpar'}\bm E_0(z',\bm k_i)$, so all
$\rpar'$ dependence sits in three exponentials and the lateral integral is
elementary:
\begin{equation}
  \int\dd^2r_\parallel'\;
  \ee^{\ii\bm q\cdot\rpar'}\,
  \ee^{-\ii\bm h\cdot\rpar'}\,
  \ee^{-\ii\bm k_{i\parallel}\cdot\rpar'}
  =(2\pi)^2\,\delta^{(2)}\!\left(\bm q-\bm k_{i\parallel}-\bm h\right).
  \label{eq:lateral-delta}
\end{equation}
The delta function is the Laue condition: for each rod exactly one
angular-spectrum mode survives,
\begin{equation}
  \bm q=\bm k_{i\parallel}+\bm h\equiv\bm k_{f\parallel},
  \qquad\text{that is}\qquad
  \Qpar=\bm h,
  \label{eq:laue}
\end{equation}
at any incidence angle, any exit angle, and any azimuth.  The $1/r$ spherical
falloff has cancelled against the density of angular-spectrum modes: what
emerges is a genuine, non-decaying plane wave per rod.

\subsection{The matrix element and the rod amplitude}

Collecting the constants of Eqs.~\eqref{eq:dwba-master}, \eqref{eq:separation},
\eqref{eq:green-1d} and \eqref{eq:lateral-delta},
\begin{equation}
  -\kzero^2\cdot\frac{1}{(2\pi)^2}
  \cdot\left(-\frac{\ii}{2k_{z,0,f}}\right)\cdot(2\pi)^2
  =+\frac{\ii\kzero^2}{2k_{z,0,f}},
  \label{eq:constants}
\end{equation}
and converting from $\delta n^2$ to $\Delta\rho_f$ with
Eq.~\eqref{eq:density-contrast}, the $\kzero^2$ cancels exactly.  Defining the
matrix element of rod $\bm h$, normalised to one reference lateral cell,
\begin{equation}
  \boxed{
  F_{\bm h}
  =\int\dd z\int_{\rm cell}\dd^2r_\parallel\;
   \Delta\rho_f(\rpar,z)\,
   \ee^{\ii\bm h\cdot\rpar}\,
   \bm E_0(z,-\bm k_f)
   \mathbin{\cdot}
   \bm E_0(z,\bm k_i),
  }
  \label{eq:Fh-definition}
\end{equation}
the scattered field on rod $\bm h$ is the plane wave
\begin{equation}
  \boxed{
  r_{\bm h}
  =-\frac{2\pi\ii r_e}{A_{\rm ref}\,k_{z,0,f}}\,F_{\bm h}.
  }
  \label{eq:rod-amplitude}
\end{equation}

Equation~\eqref{eq:Fh-definition} is the only matrix element in this document.
It is used unchanged for every rod:
\begin{itemize}
  \item for $\bm h\neq\bm 0$ its integrand contains no reference term, by
        Eq.~\eqref{eq:coeff-nonspec};
  \item for $\bm h=\bm 0$ its integrand contains the constant
        $\delta_m+\ii\beta_m$ density of every prepared slab, by
        Eq.~\eqref{eq:coeff-spec}, through $\Delta\rho_f$'s
        $-\sum_m\overline\rho_{f,m}\mathbf 1_{I_m}$ term.
\end{itemize}
Chapter~\ref{chap:unit-cells} evaluates Eq.~\eqref{eq:Fh-definition} for atoms,
producing the four coherent channels.  Chapter~\ref{chap:specular} evaluates
the reference part of $F_{\bm 0}$ and partitions it among generated records.
The reciprocal exit field in Eq.~\eqref{eq:Fh-definition} is not complex
conjugated; conjugation enters only after the coherent amplitude has been
formed.

\subsection{Born--Fresnel check}
\label{sec:born-check}

Take a uniform half-space $z<0$ of density $\rho_f$ against a vacuum
reference, so that only $\bm h=\bm 0$ contributes, the field product reduces
to $\ee^{\ii Q_zz}$, and
$\int_{-\infty}^{0}\ee^{\ii Q_zz}\dd z=1/(\ii Q_z)$.  Then
$F_{\bm 0}=A_{\rm ref}\rho_f/(\ii Q_z)$ and, with $Q_z=-2k_{z,0}>0$ from
Eq.~\eqref{eq:channel-qz},
\begin{equation}
  r_{\bm 0}
  =-\frac{2\pi\ii r_e}{A_{\rm ref}k_{z,0}}\cdot
    \frac{A_{\rm ref}\rho_f}{\ii Q_z}
  =-\frac{2\pi r_e\rho_f}{k_{z,0}Q_z}
  =+\frac{4\pi r_e\rho_f}{Q_z^{2}}\;>0,
  \label{eq:born-fresnel}
\end{equation}
which is the known kinematic reflectivity amplitude, with the positive sign of
the exact Fresnel expansion $r=(k_{z,0}-k_{z,1})/(k_{z,0}+k_{z,1})$ for $n<1$.
Both the sign and the constant in Eq.~\eqref{eq:rod-amplitude} are therefore
fixed independently of any external sign convention.

\chapter{Evaluating \texorpdfstring{$F_{\bm h}$}{Fh}: atoms and the four coherent channels}
\label{chap:unit-cells}

This chapter evaluates the matrix element $F_{\bm h}$ of
Eq.~\eqref{eq:Fh-definition} for the atomic part of the contrast.  The
reference part, which exists only at $\bm h=\bm0$, is evaluated in
Chapter~\ref{chap:specular}.  Both are contributions to the same
$F_{\bm h}$; they are never separate matrix elements.

\section{Where the four channels come from}

Section~\ref{sec:green} showed that the reciprocal exit field entering
Eq.~\eqref{eq:Fh-definition} is $u_-$, the solution of the homogeneous
stratified problem selected by the boundary condition below the stack.  Inside
prepared medium $m$ that solution has two branches, and so does the incident
field:
\begin{equation}
  \bm E_0(z',-\bm k_f)
  =A_{m,f}^{+}\ee^{-\ii k_{z,m,f}z'}+A_{m,f}^{-}\ee^{+\ii k_{z,m,f}z'},
  \qquad
  \bm E_0(z',\bm k_i)
  =A_{m,i}^{+}\ee^{-\ii k_{z,m,i}z'}+A_{m,i}^{-}\ee^{+\ii k_{z,m,i}z'}.
  \label{eq:two-branches}
\end{equation}
Their product in Eq.~\eqref{eq:Fh-definition} therefore contains four terms,
$A_{m,i}^{\sigma_i}A_{m,f}^{\sigma_f}$.  The four channels are not an
additional model: they are the two branches of each reference field, and
nothing more needs to be assumed to justify them.  Let
\begin{equation}
  \sigma_i,\sigma_f\in\{+1,-1\},
  \qquad
  \nu=(\sigma_i,\sigma_f),
\end{equation}
with the fixed channel order
\begin{equation}
  ++,\quad +-,\quad -+,\quad --.
  \label{eq:channel-order}
\end{equation}

\section{Channel wavevectors and coefficients}

For medium $m$, channel $\nu$ has
\begin{equation}
  \bm Q_{m}^{\nu}
  =\left(\Qpar,Q_{z,m}^{\nu}\right),
  \qquad
  Q_{z,m}^{\nu}
  =-\left(\sigma_i k_{z,m,i}+\sigma_f k_{z,m,f}\right),
  \label{eq:channel-qz}
\end{equation}
with $\Qpar=\bm h$ by the Laue condition, Eq.~\eqref{eq:laue}, and
\begin{equation}
  \left(Q_m^\nu\right)^2
  =|\bm h|^2+\left(Q_{z,m}^{\nu}\right)^2.
  \label{eq:channel-q2}
\end{equation}
Neither $Q_{z,m}^{\nu}$ nor $(Q_m^\nu)^2$ is generally real in an absorbing
medium.  Because Eq.~\eqref{eq:renaud-kz} makes $\operatorname{Re}k_{z,m}\le0$,
Eq.~\eqref{eq:channel-qz} gives $Q_z>0$ for the $++$ channel, as it must.

Let $A_{m,i}^{\sigma_i}$ and $A_{m,f}^{\sigma_f}$ be the local Renaud
electric-field amplitudes of Eq.~\eqref{eq:two-branches}, referenced by the
field solver to $z_{m,i}^{\rm ref}$ and $z_{m,f}^{\rm ref}$.  The complete
coefficient multiplying a Fourier factor $\exp(\ii\bm Q_m^\nu\cdot\bm r)$ is
\begin{align}
  \Cchan_{m,\nu}
  ={}&A_{m,i}^{\sigma_i}A_{m,f}^{\sigma_f}
  P_{m,\nu}\nonumber\\
  &\times\exp\!\left\{
    \ii\left[
      \sigma_i k_{z,m,i}z_{m,i}^{\rm ref}
      +\sigma_f k_{z,m,f}z_{m,f}^{\rm ref}
    \right]
  \right\}.
  \label{eq:channel-coefficient}
\end{align}
The reference phase in Eq.~\eqref{eq:channel-coefficient} is not optional: the
stored $A^\pm$ are local amplitudes, not amplitudes referenced to the global
crystallographic origin.

For the right-handed Renaud polarization basis used by orGUI, the bilinear
polarization contractions are
\begin{align}
  P_{ss}^{\sigma_i\sigma_f}
    &=\cos\varphi,\nonumber\\
  P_{sp}^{\sigma_i\sigma_f}
    &=-\sigma_f\sin\alpha_f\sin\varphi,\nonumber\\
  P_{ps}^{\sigma_i\sigma_f}
    &=-\sigma_i\sin\alpha_i\sin\varphi,\nonumber\\
  P_{pp}^{\sigma_i\sigma_f}
    &=\sin\alpha_i\sin\alpha_f
      \left[
        \frac{\kzero^2\cos\alpha_i\cos\alpha_f}
             {k_{z,m,i}k_{z,m,f}}
        -\sigma_i\sigma_f\cos\varphi
      \right].
  \label{eq:polarization-contractions}
\end{align}
These are evaluated at the nominal $\bm k_i$ and $\bm k_f$; by
Sec.~\ref{sec:why-not-farfield} that remains accurate to $\sim L/R$ even where
a far-field treatment of the \emph{phase} would fail badly.

The solver's upward internal tangential coefficient for $p$ polarization has
the opposite sign from the public Renaud amplitude.  Thus
\begin{equation}
  A_{p,m}^{+}=T_{p,m}^{+},
  \qquad
  A_{p,m}^{-}=-T_{p,m}^{-}.
  \label{eq:p-sign}
\end{equation}
Implementations should use the solver's analytic $A^\sigma/k_z$ limits in the
$pp$ expression rather than divide numerically at a critical angle.

\section{One generated UnitCell record}

The surface model is flattened into generated unit-cell records $u$.  Before
flattening, every source \texttt{UnitCell} is split into the same ordered
structural-layer views used to build the optical profile.  A record therefore
contains one such layer cell, a crystal-model weight $w_u$, its coherent domain
transforms $d$, and possibly signed domain occupancies $p_{ud}$.  A one-layer
source cell is already a single record.  Let $A_u$ be the source cell's lateral
area and $A_{\rm ref}$ the crystal-wide reference area of
Eq.~\eqref{eq:lattice-expansion}.  The required area conversion is
\begin{equation}
  s_u=\frac{A_{\rm ref}}{A_u}.
  \label{eq:area-scale}
\end{equation}
Current Film and interface domain transforms preserve lateral area.  If a
future model permits in-plane strain, $A_u$ in Eq.~\eqref{eq:area-scale} must be
replaced by the transformed projected area, including its Jacobian.

After stacking and the coherent-domain transform, atom $a$ of record $u$ and
domain $d$ is at the physical Cartesian position $\bm r_{uad}$.  The atom is
assigned to the unique prepared optical interval
\begin{equation}
  m(u,a,d)\quad\hbox{such that}\quad z_{uad}\in I_{m(u,a,d)}.
  \label{eq:medium-assignment}
\end{equation}
A generated unit cell that crosses an optical boundary must therefore be split
by atom.  Assigning one medium to the entire cell would use the wrong $A^\pm$,
$k_z$, polarization contraction, and complex form-factor argument for some of
its atoms.

For channel $\nu$ in medium $m$, define the atomic contribution
\begin{align}
  \mathcal A_{u,m}^{\nu}
  ={}&s_u w_u\Cchan_{m,\nu}
  \sum_d p_{ud}
  \sum_{\substack{a\in u\\m(u,a,d)=m}}
  o_a f_a\!\left((Q_m^\nu)^2,E\right)
  \exp\!\left[-W_a\!\left(\bm Q_m^\nu\right)\right]
  \exp\!\left(\ii\bm Q_m^\nu\mathbin{\cdot}\bm r_{uad}\right).
  \label{eq:atomic-channel}
\end{align}
Here $o_a$ is the atomic occupancy.  For the currently implemented in-plane
and out-of-plane Debye--Waller parameters,
\begin{equation}
  W_a(\bm Q_m^\nu)
  =\frac{B_{a,\parallel}|\bm h|^2
          +B_{a,\perp}(Q_{z,m}^\nu)^2}{16\pi^2}.
  \label{eq:debye-waller}
\end{equation}
Equations~\eqref{eq:atomic-factor}, \eqref{eq:channel-q2}, and
\eqref{eq:debye-waller} are evaluated at complex arguments without taking an
absolute value or a real part.

\section{The atomic contribution to \texorpdfstring{$F_{\bm h}$}{Fh}}

The atomic part of record $u$'s contribution to Eq.~\eqref{eq:Fh-definition} is
\begin{equation}
  \boxed{
  F_{\bm h,u}^{\rm at}
  =\sum_m\sum_{\nu\in\{++,+-,-+,--\}}
   \mathcal A_{u,m}^{\nu},
  }
  \label{eq:Fh-atomic}
\end{equation}
and for $\bm h\neq\bm 0$ this is the whole of it,
\begin{equation}
  F_{\bm h}=\sum_u F_{\bm h,u}^{\rm at},
  \qquad \bm h\neq\bm 0,
  \label{eq:Fh-nonspecular}
\end{equation}
because Eq.~\eqref{eq:coeff-nonspec} carries no reference term.  The specular
coefficient $F_{\bm 0}$ carries one additional contribution and is completed in
Chapter~\ref{chap:specular}.

The four terms in Eq.~\eqref{eq:Fh-atomic} are amplitudes and must be added
before an intensity is formed.  The incorrect expression
$\sum_\nu|\mathcal A_{u,m}^{\nu}|^2$ discards the standing-wave interference
between incident and exit branches that Eq.~\eqref{eq:two-branches} makes
explicit.

\section{Useful limits}

In a terminal substrate, the boundary condition that defines $u_-$ in
Eq.~\eqref{eq:bc-minus} gives
\begin{equation}
  A_{M,i}^{-}=A_{M,f}^{-}=0,
  \label{eq:terminal-branches}
\end{equation}
so only the $++$ channel survives.  In the high-angle Born limit of a finite
cell in vacuum, $A_i^+=A_f^+=1$ and all reflected branches vanish; the same
reduction occurs.  These are limiting cases of the four-channel expression,
not reasons to omit channels in finite graded layers.

If the cell is wholly contained in one medium, Eq.~\eqref{eq:atomic-channel}
reduces to four ordinary structure-factor evaluations at the four complex
values $Q_{z,m}^\nu$.  If it crosses a boundary, there is no single structure
factor that can be multiplied by one set of channel coefficients; the
atom-medium partition in Eq.~\eqref{eq:atomic-channel} is essential.

\chapter{The specular coefficient \texorpdfstring{$F_{\bm 0}$}{F0}}
\label{chap:specular}

\section{What is different at \texorpdfstring{$\bm h=\bm 0$}{h = 0}}

Everything that distinguishes the specular rod was already established in
Eqs.~\eqref{eq:coeff-nonspec}--\eqref{eq:coeff-spec}: the prepared reference
$n_0^2(z)$ is laterally uniform, so it appears in the $\bm h=\bm 0$ Fourier
coefficient of $\delta n^2$ and in no other.  Equivalently, in the language of
Eq.~\eqref{eq:lateral-delta}, a laterally uniform density has all of its
lateral Fourier weight at $\bm q-\bm k_{i\parallel}=\bm 0$, and a genuine rod
$\bm h\neq\bm 0$ never sits there.

Consequently $F_{\bm 0}$, and only $F_{\bm 0}$, receives a second contribution
besides the atomic sum of Eq.~\eqref{eq:Fh-atomic}:
\begin{equation}
  F_{\bm 0}=\sum_u\left(F_{\bm 0,u}^{\rm at}+F_{\bm 0,u}^{\rm ref}\right),
  \qquad
  F_{\bm h}=\sum_uF_{\bm h,u}^{\rm at}\ \ (\bm h\neq\bm 0).
  \label{eq:Fh-two-parts}
\end{equation}
This is one matrix element with two kinds of contribution, not two matrix
elements.  In particular nothing is ``subtracted afterwards'': the sign is
already carried by $\Delta\rho_f$ in Eq.~\eqref{eq:density-contrast}.

It is worth stating what would go wrong if the reference contribution were
instead applied cell by cell off-specular.  Doing so would replace the exact
lateral delta function of Eq.~\eqref{eq:lateral-delta} by a finite-cell shape
transform and manufacture diffuse scattering that the reference profile does
not produce.  Instrumental resolution may mix genuine diffuse intensity into a
detector pixel on the specular rod, but that belongs to a later resolution
convolution, not to an interpolation of Eq.~\eqref{eq:Fh-two-parts} near
$\bm h=\bm 0$.

\section{Integral of one constant optical slab}

Within prepared medium $m$, both $\overline\rho_{f,m}$ and every channel
coefficient are constant; only the propagation phase varies with $z$.  Define
the stable depth integral
\begin{equation}
  \Phiint(q;a,b)
  =\int_a^b\ee^{\ii qz}\,\dd z
  =\begin{cases}
    \dfrac{\ee^{\ii qb}-\ee^{\ii qa}}{\ii q}, & q\neq0,\\[3mm]
    b-a, & q=0.
  \end{cases}
  \label{eq:phi-finite}
\end{equation}
For numerical work, an equivalent midpoint form is preferable:
\begin{equation}
  \Phiint(q;a,b)
  =(b-a)\,
   \ee^{\ii q(a+b)/2}
   \operatorname{sinc}\!\left(\frac{q(b-a)}{2}\right),
  \qquad
  \operatorname{sinc}x=\frac{\sin x}{x},
  \label{eq:phi-sinc}
\end{equation}
or the numerator of Eq.~\eqref{eq:phi-finite} may be evaluated with a complex
\texttt{expm1}.  For a lower terminal half-space, with the physical decaying
branch,
\begin{equation}
  \Phiint(q;-\infty,b)=\frac{\ee^{\ii qb}}{\ii q}.
  \label{eq:phi-half-space}
\end{equation}

Carrying the $-\sum_m\overline\rho_{f,m}\mathbf 1_{I_m}$ term of
Eq.~\eqref{eq:density-contrast} through Eq.~\eqref{eq:Fh-definition}, the
reference contribution to $F_{\bm 0}$ at profile level is
\begin{equation}
  \boxed{
  \sum_uF_{\bm 0,u}^{\rm ref}
  =-A_{\rm ref}
   \sum_m\overline\rho_{f,m}
   \sum_{\nu\in\{++,+-,-+,--\}}
   \Cchan_{m,\nu}
   \Phiint\!\left(Q_{z,m}^{\nu};z_m^-,z_m^+\right).
  }
  \label{eq:profile-reference-term}
\end{equation}
Finite graded layers generally retain all four channels in
Eq.~\eqref{eq:profile-reference-term}; only the terminal substrate reduces to
$++$, by Eq.~\eqref{eq:terminal-branches}.

\section{A well-defined per-record partition}

Only the total reference profile is physical.  A decomposition of that profile
among overlapping model objects is bookkeeping and is not unique unless a
rule is fixed.  orGUI can make the decomposition deterministic by retaining
the forward-density contribution of every generated record while building the
profile.

Let $\overline\rho_{f,u,m}$ denote record $u$'s share of the constant reference
density in prepared medium $m$.  The required conservation rule is
\begin{equation}
  \boxed{
  \sum_u\overline\rho_{f,u,m}=\overline\rho_{f,m}
  \quad\text{for every prepared medium }m.
  }
  \label{eq:reference-partition}
\end{equation}
Before optical-profile simplification, a generated slab of physical thickness
$t_{ud}$ contributes naturally
\begin{equation}
  \overline\rho_{f,ud}
  =\frac{w_u p_{ud}}{A_u t_{ud}}
   \sum_{a\in u}o_a f_a(0,E).
  \label{eq:natural-density-share}
\end{equation}
Normal strain changes $t_{ud}$ and hence the volume density, but preserves the
integrated areal content.  Signed occupancies in interface and etching models
remain signed in Eq.~\eqref{eq:natural-density-share}.

Coincident profile contributions are added.  If adjacent optical media are
merged, first determine the merged boundaries from the total profile, then
thickness-average every record's density over those same boundaries.  This is
a linear operation once the grouping is fixed and preserves
Eq.~\eqref{eq:reference-partition}.  Reconstructing shares independently with
different merge decisions would not preserve the prepared reference.

With Eq.~\eqref{eq:reference-partition} fixed, record $u$'s reference
contribution is
\begin{equation}
  F_{\bm 0,u}^{\rm ref}
  =-A_{\rm ref}
   \sum_m\overline\rho_{f,u,m}
   \sum_\nu\Cchan_{m,\nu}
   \Phiint\!\left(Q_{z,m}^{\nu};z_m^-,z_m^+\right),
  \label{eq:per-cell-reference}
\end{equation}
so that record $u$'s complete share of the specular matrix element is
\begin{equation}
  \boxed{
  F_{\bm 0,u}
  =\sum_m\sum_\nu\mathcal A_{u,m}^{\nu}
   -A_{\rm ref}
    \sum_m\overline\rho_{f,u,m}
    \sum_\nu\Cchan_{m,\nu}
    \Phiint\!\left(Q_{z,m}^{\nu};z_m^-,z_m^+\right).
  }
  \label{eq:per-cell-specular}
\end{equation}
Summing Eq.~\eqref{eq:per-cell-specular} over $u$ and applying
Eq.~\eqref{eq:reference-partition} recovers
Eq.~\eqref{eq:profile-reference-term} and hence the profile-level
$F_{\bm 0}$ of Eq.~\eqref{eq:Fh-two-parts}:
\begin{equation}
  \boxed{
  F_{\bm 0}
  =\sum_uF_{\bm 0,u}
  =\sum_uF_{\bm 0,u}^{\rm at}
   -A_{\rm ref}
    \sum_m\overline\rho_{f,m}
    \sum_\nu\Cchan_{m,\nu}
    \Phiint\!\left(Q_{z,m}^{\nu};z_m^-,z_m^+\right).
  }
  \label{eq:total-specular}
\end{equation}
If callers do not need an attribution by object,
Eq.~\eqref{eq:total-specular} is the preferred calculation: accumulate all
atomic terms, evaluate the prepared reference once, and add once.

\section{Reduction to the implemented terminal bulk formula}

Let the top terminal bulk cell occupy $b-c\leq z<b$, with repeat vector
$-\bm c$ into the substrate and $c=c_z>0$.  Only the $++$ channel survives, by
Eq.~\eqref{eq:terminal-branches}.  If $F_{\rm uc}$ denotes the ordinary
unit-cell structure factor including the global phase of the top cell, the
semi-infinite atomic contribution to $F_{\bm 0}$ is
\begin{equation}
  F_{\bm 0,\rm bulk}^{\rm at}
  =\Cchan_{M,++}
   \frac{F_{\rm uc}}
        {1-\exp(-\ii\bm Q_M^{++}\mathbin{\cdot}\bm c)},
  \label{eq:bulk-atomic-sum}
\end{equation}
and, for a surface-normal repeat, the mean-density contribution is
\begin{equation}
  F_{\bm 0,\rm bulk}^{\rm ref}
  =-\Cchan_{M,++}A_{\rm ref}\overline\rho_{f,M}
   \frac{\ee^{\ii Q_{z,M}^{++}b}}{\ii Q_{z,M}^{++}}.
  \label{eq:bulk-reference}
\end{equation}
Their sum is the previously implemented bulk result.  The four-channel
finite-layer formula therefore extends the existing kernel without changing
the terminal half-space convention.

Equation~\eqref{eq:bulk-reference} also shows why the reference contribution
cannot be dropped: it is precisely the $Q_z\to0$ divergence of the kinematic
lattice sum in Eq.~\eqref{eq:bulk-atomic-sum}, that is, the $\bm h=\bm 0$
component that the prepared optical profile already describes and that $r_0$
in Eq.~\eqref{eq:specular-reflectivity} already carries.  Retaining both
without Eq.~\eqref{eq:total-specular} would count it twice.

\chapter{Mapping to orGUI models and validation}
\label{chap:mapping}

\section{Flattened atomic representation}

The mathematical kernel should consume a flattened list of generated atomic
records rather than special-case each high-level model.  After stacking, the
records have the following meanings.

\begin{longtable}{@{}p{0.23\textwidth}p{0.67\textwidth}@{}}
\toprule
Model & Records supplied to the $F_{\bm h}$ sum \\
\midrule
\endfirsthead
\toprule
Model & Records supplied to the $F_{\bm h}$ sum \\
\midrule
\endhead
\texttt{UnitCell}
& Ordered structural-layer views, atoms, source lateral area, internal
coherent-domain transforms and occupancies, and the crystal-wide component
weight. \\
\texttt{Film}
& The finite cyclic layer cells generated by \texttt{createLayers}.  Their
domain translations place them at the physical stacked heights. \\
\texttt{EpitaxyInterface}
& Positive upper-material probability records; positive lower-material
records at their possibly strain-coupled positions; and negative lower-
material records that remove the sharp, unstrained bulk already present in
the terminal lattice sum. \\
\texttt{PoissonSurface}
& Positive records for grown material and negative records for material
removed from the sharp Film termination.  Termination corrections remain
coherent with the Film amplitude. \\
\bottomrule
\end{longtable}

Negative occupancies are amplitudes, not negative intensities.  They must pass
unchanged through Eqs.~\eqref{eq:atomic-channel} and
\eqref{eq:natural-density-share}.  Likewise, alternative crystal domains are
applied as coherent amplitudes within each configured domain; intensities are
combined only at the model level where the existing API explicitly defines an
incoherent mixture.

Water and other continuous-density models do not have an atomic UnitCell sum
and are outside this chapter.  They may use the same slab integral in
Eq.~\eqref{eq:profile-reference-term}, but their actual-density integral must
be derived from their own continuous profile.

\section{Preparation and evaluation boundary}

The prepared object must freeze all quantities that define the reference
problem:
\begin{itemize}
  \item the physical optical boundaries and the final, possibly simplified,
        values $\delta_m,\beta_m$;
  \item incident and reciprocal exit $k_z$, $A^\pm$, $A^\pm/k_z$, and reference
        planes for every medium;
  \item incident and exit polarization, external geometry, and the rod
        $\bm h=\Qpar$;
  \item the reference-area convention and the per-record reference shares
        satisfying Eq.~\eqref{eq:reference-partition}.
\end{itemize}
Atomic positions, occupancies, Debye--Waller parameters, and species may be
reevaluated against a frozen reference if that is the intended fit operation.
Changing lattice transforms, stacking geometry, photon energy, optical
boundaries, profile simplification, or the reference area invalidates the
preparation.  Rebuilding only the atomic list while silently changing the
optical profile would combine two different reference problems.

For each scattering point, evaluation is naturally ordered as follows:
\begin{enumerate}
  \item identify the rod from the Laue condition, Eq.~\eqref{eq:laue}, and
        calculate $\bm Q_m^\nu$, $P_{m,\nu}$, and $\Cchan_{m,\nu}$ for all
        media and all four channels;
  \item assign every transformed atom to its physical optical medium;
  \item accumulate Eq.~\eqref{eq:atomic-channel} coherently over atoms,
        domains, media, and channels, giving $F_{\bm h,u}^{\rm at}$;
  \item if $\bm h=\bm 0$, add the constant-profile contribution of
        Eq.~\eqref{eq:profile-reference-term}, completing $F_{\bm 0}$;
        otherwise $F_{\bm h}$ is already complete;
  \item convert to a rod amplitude with Eq.~\eqref{eq:rod-amplitude}, and only
        then form an intensity or combine with the prepared unperturbed
        reflection amplitude.
\end{enumerate}

\section{Observables}

Equation~\eqref{eq:rod-amplitude} already gives the amplitude of the emergent
plane wave on rod $\bm h$,
\[
  r_{\bm h}=-\frac{2\pi\ii r_e}{A_{\rm ref}k_{z,0,f}}F_{\bm h},
\]
normalised to one reference lateral cell, for every rod and every scattering
geometry.  Nothing further needs to be assumed about the measurement geometry:
fixed $\alpha_i$, fixed $\alpha_f$, $\alpha_i=\alpha_f$, coplanar or not, the
scan mode only decides which $(\alpha_i,\alpha_f)$ pair --- and hence which
$k_{z,m}$ and which $A^\pm$ --- are evaluated at each measured point.

On a genuine crystal truncation rod, $\bm h\neq\bm 0$, there is no unperturbed
beam in that direction and the observable follows from $r_{\bm h}$ alone.  On
the specular rod the scattered field is combined coherently with the prepared
plane-wave reflection amplitude,
\begin{equation}
  r_{\rm total}=r_0+r_{\bm 0},
  \qquad
  R_{\rm coh}=|r_{\rm total}|^2.
  \label{eq:specular-reflectivity}
\end{equation}
The $r_0$ in Eq.~\eqref{eq:specular-reflectivity} must be the same $r_0$ that
appears in Eq.~\eqref{eq:u-minus}, that is, it must come from the same prepared
incident field as the $A_{m,i}^{\pm}$ used in the matrix element.  For
unpolarized radiation, calculate $s$ and $p$ independently and average their
intensities; do not add their amplitudes.

\paragraph{Field intensity and the experimental input.}
The fitting design uses $R=|r|^2$ as a dimensionless field-intensity ratio,
including off-specular rod amplitudes.  Experimental data must already be
corrected to this convention and absolutely normalized when the fitted scale
is fixed to one.  For plane waves in vacuum, the ratio of normal powers
through a surface parallel to the sample is instead
$(\kappa_f/\kappa_i)|r|^2$, where
$\kappa_{i,f}=\kzero\sin\alpha_{i,f}>0$.  These two observables coincide in
specular reflection.  The normal-power ratio is not a universal conversion
from detector counts: illuminated area, footprint, and detector acceptance
belong to experimental reduction.  The fitter therefore adds no such ratio
and does not redefine $r$.  This follows the distinction between field
intensity, flux conservation, and differential scattering cross sections in
Renaud et al., Secs.~5.2.6 and 5.3.3--5.3.6.  Within each polarization pair,
the four DWBA propagation branches remain a coherent amplitude sum;
incoherent polarization mixtures are combined only after squaring.

\section{Returned quantities}

The kernel returns one matrix element and everything else is derived from it
by the prefactor of Eq.~\eqref{eq:rod-amplitude}.  Implementations write that
prefactor with the positive vacuum exit wavevector
$\kappa_f=\kzero\sin\alpha_f=-k_{z,0,f}$, so the sign appears reversed relative
to Eq.~\eqref{eq:rod-amplitude}; the two are identical.

\begin{center}
\begin{tabular}{@{}lll@{}}
\toprule
Quantity & Definition & Units \\
\midrule
$F_{\bm h}$
  & Eq.~\eqref{eq:Fh-definition}
  & electrons per $A_{\rm ref}$ \\
$r_{\bm h}$
  & $2\pi\ii r_e F_{\bm h}/(A_{\rm ref}\kappa_f)$
  & dimensionless \\
$r_0$
  & Fresnel amplitude of the prepared reference, Eq.~\eqref{eq:u-minus}
  & dimensionless \\
$r$
  & $r_0+r_{\bm h}$, with $r_0=0$ for $\bm h\neq\bm 0$
  & dimensionless \\
$F_{\rm eff}$
  & $A_{\rm ref}\kappa_f\,r/(2\pi\ii r_e)$
  & electrons per $A_{\rm ref}$ \\
$R$
  & $|r|^2$
  & dimensionless \\
\bottomrule
\end{tabular}
\end{center}

Three points are worth stating explicitly, because each has caused confusion.

\paragraph{Where $r_e$ belongs.}  The classical electron radius enters once,
inside the prefactor.  The separate convention, in which a bare matrix element
is turned into a cross-section kernel by $r_e^2|F_{\bm h}|^2$, applies to
$F_{\bm h}$ and not to $r$.  The two must never be mixed.

\paragraph{What $F_{\rm eff}$ is.}  Inverting the prefactor on the
\emph{total} amplitude gives the structure factor that a kinematic analysis
would infer from this reflectivity.  It is a DWBA output; it is not the
structure factor of a separate kinematical model, which the tutorial material
writes $F_{\rm kin}$.  The two are compared against each other, never
substituted for one another.  It is not linear in the model density,
because $r_0$ is zeroth order and exact in the reference to all orders, so it
is not the DWBA structure factor of the sample; that is $F_{\bm h}$.  It
vanishes as $\alpha_f\to0$ and is meaningful only well above the critical
angle.  Off-specular, where $r_0=0$, it degenerates to $F_{\bm h}$.  Its value
is the large-$Q_z$ limit, where $r_0\to4\pi r_e\overline\rho_f/Q_z^2$ and
$F_{\rm eff}\to F_{\bm h}+A_{\rm ref}\overline\rho_f/(\ii Q_z)$, that is, the
full kinematic structure factor of the actual density --- exactly the quantity
a kinematically reduced data set contains.

\paragraph{Polarization-reduced fitting magnitudes.}
Reserve $P$ for the conventional dimensionless polarization intensity factor,
and denote the amplitude prefactor by
$c=2\pi\ii r_e/(A_{\rm ref}\kappa_f)$.  If experimental structure-factor
data were polarization corrected, intensity was multiplied by $1/P$ and
amplitude magnitude by $1/\sqrt{P}$.  Public fitting metadata stores $P$;
an absent factor means no correction, equivalently $P=1$.  Let $\overline R$
be the resolution-broadened field intensity, with any incident mixture and
unanalysed outgoing states summed incoherently.  The predicted magnitude is
\[
  F_{\rm pred}=\frac{\sqrt{\overline R/P}}{|c|},
\]
using central-geometry $P$ and $c$.  There is no extra $\kappa_i$ in this
conversion, although the incident geometry still enters the wavefields.
For a narrow integrated region where $P$ changes little, a central factor
also approximates pixelwise polarization correction.  This is not an exact
identity for general integration weights; it must be adequate for the
experimental uncertainty.  It does not replace the polarization-dependent
DWBA fields or imply a unique complex $F_{\rm eff}$ for a mixed measurement.

\paragraph{Why $R=|r|^2$ is the reflectivity to use.}  On the specular rod at
small angles the problem is genuinely one-dimensional, so the scalar reduction
of Eq.~\eqref{eq:helmholtz} involves no approximation and the $s$ and $p$
channels become degenerate; and $r_0$ carries the exact Fresnel and refraction
response of the prepared reference to all orders.  The only error is therefore
second order in the density contrast, and it shrinks with decreasing angle:
below the critical angle the exact reflectivity is unity and $|r_0+r_{\bm h}|^2$
reproduces it.  A reference whose density is wrong by $10\%$ gives a relative
error near $10^{-4}$ at one tenth of the critical angle, rising to a few
percent near $\alpha_c$, and falling quadratically as the contrast error is
reduced.  Squaring the amplitude rather than expanding it also keeps $R$
non-negative.  The strictly linearised variant,
$|r_0|^2+2\operatorname{Re}(r_0^*r_{\bm h})$, is useful only as a diagnostic:
a large difference between the two means $|r_{\bm h}|$ is not small against
$|r_0|$ and the first-order treatment is being pushed beyond its regime.

\section{Regression invariants}

The following checks are required before the atomistic extension is considered
scientifically usable.

\begin{enumerate}
  \item \textbf{Four-channel coherence.}  A finite layer with nonzero
        $A_i^\pm$ and $A_f^\pm$ equals the direct atom-by-atom evaluation of
        Eq.~\eqref{eq:Fh-definition}; changing to a sum of channel intensities
        must fail the test.
  \item \textbf{Born limit.}  Setting $A_i^+=A_f^+=1$ and the other branches to
        zero recovers the conventional UnitCell structure factor and the
        standard polarization contraction.
  \item \textbf{Born--Fresnel amplitude.}  A uniform half-space against a
        vacuum reference reproduces Eq.~\eqref{eq:born-fresnel}, including its
        sign.  This pins the prefactor of Eq.~\eqref{eq:rod-amplitude}
        independently of any external convention.
  \item \textbf{Terminal limit.}  A finite stack made progressively thicker
        converges to Eqs.~\eqref{eq:bulk-atomic-sum} and
        \eqref{eq:bulk-reference} away from exact Bragg poles.
  \item \textbf{Identical actual and reference density.}  Replacing every
        optical slab by its exactly matching atomic forward density gives
        $F_{\bm 0}=0$ within the atomic-truncation accuracy.
  \item \textbf{Off-specular invariance.}  Adding the reference-density
        bookkeeping does not change any $F_{\bm h}$ with $\bm h\neq\bm0$, as
        required by Eq.~\eqref{eq:coeff-nonspec}.
  \item \textbf{Partition invariance.}  Summing per-record results from
        Eq.~\eqref{eq:per-cell-specular} equals the profile-level
        Eq.~\eqref{eq:total-specular}, including after profile coalescing and
        simplification.
  \item \textbf{Area invariance.}  Two commensurate descriptions with different
        source lateral cells give the same amplitude per $A_{\rm ref}$ after
        Eq.~\eqref{eq:area-scale} is applied to both atomic and continuum
        terms.
  \item \textbf{Signed-model identity.}  At zero width, zero growth, or zero
        etching, the positive and negative interface/surface records cancel in
        both contributions to $F_{\bm h}$.
  \item \textbf{Boundary split.}  Moving an atom across an optical boundary
        changes its channel coefficient and $Q_z$ exactly once; a cell
        straddling a boundary agrees with the same atom list represented as two
        records.
  \item \textbf{Low-$Q_z$ stability.}  Equation~\eqref{eq:phi-finite} approaches
        the slab thickness without catastrophic cancellation.
  \item \textbf{Reciprocity.}  Exchanging source and detector geometry obeys
        the reciprocal-field symmetry of
        Eq.~\eqref{eq:u-minus-identification} without conjugating the exit
        field.
  \item \textbf{Wronskian normalisation.}  For any prepared stack, the
        Wronskian of Eq.~\eqref{eq:wronskian-value} evaluates to
        $2\ii k_{z,0}$ independently of the layer structure.  A solver whose
        $A^\pm$ violate this has an internal normalisation error that
        Eq.~\eqref{eq:rod-amplitude} would silently absorb.
  \item \textbf{No duplicate attenuation.}  Absorption carried by complex
        $Q_z$ is not also applied as an empirical depth factor.
\end{enumerate}

\section{References}

The electric-field, signed-$k_z$, reciprocal-exit-field, and four-channel
semantics follow G.~Renaud, R.~Lazzari, and F.~Leroy, ``Probing surface and
interface morphology with grazing incidence small angle X-ray scattering,''
\emph{Surface Science Reports} \textbf{64} (2009), 255--380, especially
Sections~5.2--5.4, Eq.~(91), and the four-event finite-layer examples in
Section~6.2.  The present note retains their electromagnetic derivation but
does not adopt the scalar wavefunction notation used for the optional quantum-
mechanical analogy.

Renaud's Eq.~(100) quotes the one-dimensional stratified Green function, with
its factor $1/k_z$, without derivation, citing P.~M.~Morse and H.~Feshbach,
\emph{Methods of Theoretical Physics} (McGraw--Hill, 1953) for the general
construction, together with X.~L.~Zhou and S.~H.~Chen, \emph{Physics Reports}
\textbf{257} (1995), 223--348; A.~V.~Andreev, A.~G.~Michette and A.~Renwick,
\emph{Journal of Modern Optics} \textbf{35} (1988), 1667; V.~F.~Sears,
\emph{Physical Review B} \textbf{48} (1993), 17477; I.~D.~Feranchuk,
\emph{Physical Review B} \textbf{52} (1995), 9214; and A.~A.~Maradudin and
D.~L.~Mills, \emph{Physical Review B} \textbf{11} (1975), 1392.
Section~\ref{sec:green} carries out that construction explicitly in orGUI's
own conventions, so the present document does not depend on any of them.

The terminal crystal-truncation result and the atomic-density truncation caveat
are cross-checked against V.~M.~Kaganer, ``Crystal truncation rods in
kinematical and dynamical x-ray diffraction theories,'' \emph{Physical Review
B} \textbf{75} (2007), 245425.  The perturbation as a difference from a planar
reference follows the same physical construction used by G.~H.~Vineyard,
\emph{Physical Review B} \textbf{26} (1982), 4146--4159, and by S.~K.~Sinha
\emph{et al.}, \emph{Physical Review B} \textbf{38} (1988), 2297--2311.

The companion file \texttt{dwba\_literature\_crosscheck.md} records the broader
comparison with the surface-DWBA and reflectivity literature.


\end{document}
